% BHU PYQ -- AIHC & Archaeology: Field Epigraphy (2024-25 NEP)
\documentclass[11pt,a4paper]{article}
\usepackage[a4paper,top=2.5cm,bottom=2.5cm,left=2.8cm,right=2.8cm,headheight=14pt]{geometry}
\usepackage{amsmath,amssymb}
\usepackage{fontspec}
\setmainfont{Kohinoor Devanagari}
\usepackage{microtype,fancyhdr,enumitem,hyperref,lastpage,parskip}

\hypersetup{colorlinks=true,urlcolor=black,linkcolor=black,pdftitle={AIHSE21: Field Epigraphy -- BHU NEP 2024-25}}

\pagestyle{fancy}\fancyhf{}
\renewcommand{\headrulewidth}{0.4pt}\renewcommand{\footrulewidth}{0.4pt}
\fancyhead[L]{\small \textit{Banaras Hindu University}}
\fancyhead[R]{\small \textit{AIHSE21: Field Epigraphy (2024-25)}}
\fancyfoot[C]{\small Page \thepage\ of \pageref{LastPage}}

\newlist{parts}{enumerate}{2}
\setlist[parts,1]{label=(\alph*),leftmargin=*,itemsep=4pt,topsep=3pt}
\setlist[parts,2]{label=(\roman*),leftmargin=*,itemsep=2pt,topsep=2pt}

\newcommand{\pts}[1]{\hfill \textit{[#1]}}

\begin{document}

\begin{center}
  {\large\bfseries BANARAS HINDU UNIVERSITY}\\[2pt]
  {\normalsize B. A. (Hons.) Semester II Examination, 2024-25 (NEP)}\\[6pt]
  \rule{0.8\linewidth}{0.5pt}\\[6pt]
  {\large\bfseries ANCIENT INDIAN HISTORY, CULTURE \& ARCHAEOLOGY}\\[4pt]
  {\large\bfseries Paper Code: AIHSE21 : Field Epigraphy}\\[4pt]
  {\normalsize Time: 2:30 Hours \hfill Full Marks: 70}\\[2pt]
  {\small REF NO. D-II/067}
\end{center}
\rule{\linewidth}{0.4pt}
\smallskip

\noindent \textit{Instructions: Attempt all Sections A, B and C according to instructions given. Write your Roll No.\ at the top immediately on receipt of this question paper.}

\smallskip
\noindent \textit{दिये गये निर्देशों के अनुसार खण्ड अ, ब तथा स के उत्तर दीजिये।}

\bigskip

\section*{SECTION -- A / खण्ड -- अ}
\noindent \textit{Note: Attempt any \textbf{three} questions, each in 300 words. Each question carries 10 marks.}\pts{3 $\times$ 10 = 30}

\noindent \textit{किन्हीं तीन प्रश्नों के उत्तर दीजिये, प्रत्येक उत्तर 300 शब्दों में हो। प्रत्येक प्रश्न 10 अंकों का है।}

\begin{enumerate}[label=\textbf{\arabic*.},leftmargin=*,itemsep=12pt]

  \item Define Epigraphy and discuss its importance as a primary source for reconstructing Ancient Indian History.\pts{10}
  
  अभिलेखशास्त्र की परिभाषा दीजिए तथा प्राचीन भारतीय इतिहास के पुनर्निर्माण में प्राथमिक स्रोत के रूप में इसके महत्त्व की चर्चा कीजिए।

  \item Describe the techniques and equipment used for taking estampages (paper ink squeezes) of inscriptions in the field.\pts{10}
  
  क्षेत्रीय कार्य के दौरान अभिलेखों की छाप (एस्टैम्पैज) लेने की तकनीकों तथा उपकरणों का वर्णन कीजिए।

  \item Explain the origin, development and characteristics of the Brahmi script in India.\pts{10}
  
  भारत में ब्राह्मी लिपि की उत्पत्ति, विकास एवं विशेषताओं की व्याख्या कीजिए।

  \item Discuss the decipherment of Ashokan Brahmi by James Prinsep and its historical significance.\pts{10}
  
  जेम्स प्रिंसेप द्वारा अशोकीय ब्राह्मी के वाचन (Decipherment) तथा इसके ऐतिहासिक महत्त्व की विवेचना कीजिए।

  \item Discuss Kharosthi script, its origin, direction of writing and geographical distribution.\pts{10}
  
  खरोष्ठी लिपि की उत्पत्ति, लेखन-दिशा एवं भौगोलिक वितरण की चर्चा कीजिए।

\end{enumerate}

\bigskip

\section*{SECTION -- B / खण्ड -- ब}
\noindent \textit{Note: Attempt any \textbf{four} questions, each in 150 words. Each question carries 5 marks.}\pts{4 $\times$ 5 = 20}

\noindent \textit{किन्हीं चार प्रश्नों के उत्तर दीजिये, प्रत्येक उत्तर 150 शब्दों में हो। प्रत्येक प्रश्न 5 अंकों का है।}

\begin{enumerate}[label=\textbf{\arabic*.},leftmargin=*,start=6,itemsep=10pt]

  \item Differentiate between Inscriptions on Stone (Shilalekh) and Copper-plate grants (Tamrapatra).\pts{5}
  
  शिलालेखों एवं ताम्रपत्र दानपत्रों के मध्य अन्तर स्पष्ट कीजिए।

  \item Write a short note on Prashasti as an epigraphic genre.\pts{5}
  
  अभिलेखीय विधा के रूप में 'प्रशस्ति' पर एक संक्षिप्त टिप्पणी लिखिए।

  \item Explain the role of palaeography in dating undated ancient inscriptions.\pts{5}
  
  तिथिविहीन प्राचीन अभिलेखों के काल निर्धारण में पुरालिपि शास्त्र (Palaeography) की भूमिका की व्याख्या कीजिए।

  \item Discuss the preservation methods of stone inscriptions in field sites.\pts{5}
  
  पुरास्थल क्षेत्रों में पाषाण अभिलेखों के संरक्षण एवं रख-रखाव की विधियों की चर्चा कीजिए।

  \item Write about the Allahabad Pillar Inscription (Prayag Prashasti) of Samudragupta.\pts{5}
  
  समुद्रगुप्त के प्रयाग प्रशस्ति अभिलेख के विषय में लिखिए।

  \item What is the importance of land-grant inscriptions in understanding agrarian economy?\pts{5}
  
  कृषि अर्थव्यवस्था को समझने में भूमि-दान अभिलेखों का क्या महत्त्व है?

\end{enumerate}

\bigskip

\section*{SECTION -- C / खण्ड -- स}
\noindent \textit{Note: Attempt all \textbf{ten} questions of this section in a maximum of two sentences each. Each question carries 2 marks.}\pts{10 $\times$ 2 = 20}

\noindent \textit{इस खण्ड के सभी दस प्रश्नों के उत्तर अधिकतम दो वाक्यों में दीजिये। प्रत्येक प्रश्न 2 अंकों का है।}

\begin{enumerate}[label=\textbf{\arabic*.},leftmargin=*,start=12,itemsep=8pt]
  \item 
  \begin{parts}
    \item What is Palaeography?\pts{2}
    
    पुरालिपि शास्त्र (Palaeography) क्या है?

    \item In which direction is the Kharosthi script written?\pts{2}
    
    खरोष्ठी लिपि किस दिशा में लिखी जाती है? (Right to Left / दाएँ से बाएँ)

    \item Who deciphered the Ashokan inscriptions in 1837?\pts{2}
    
    1837 ई० में अशोकीय अभिलेखों को किसने पढ़ा था? (James Prinsep)

    \item Name the language of most Ashokan edicts.\pts{2}
    
    अधिकांश अशोकीय अभिलेखों की भाषा का नाम लिखिये। (Prakrit)

    \item What is an Estampage?\pts{2}
    
    एस्टैम्पैज (छाप) क्या होता है?

    \item Name the famous inscription of King Kharavela of Kalinga.\pts{2}
    
    कलिंग के राजा खारवेल के प्रसिद्ध अभिलेख का नाम लिखिये। (Hathigumpha Inscription)

    \item Where is the Junagadh Rock Inscription of Rudradaman located?\pts{2}
    
    रुद्रदामन का जूनागढ़ शिलालेख कहाँ स्थित है? (Girnar, Gujarat)

    \item Who composed the Prayag Prashasti?\pts{2}
    
    प्रयाग प्रशस्ति की रचना किसने की थी? (Harishena)

    \item What is a Tamrapatra?\pts{2}
    
    ताम्रपत्र क्या होता है?

    \item Which era was started by King Chandragupta I in 319-320 CE?\pts{2}
    
    राजा चंद्रगुप्त प्रथम ने 319-320 ई० में किस संवत की शुरुआत की थी? (Gupta Era / गुप्त संवत)
  \end{parts}
\end{enumerate}

\end{document}
